\documentclass[12pt]{article}
\usepackage{amsmath,amssymb,booktabs,longtable,array}
\usepackage{fontspec}
\usepackage{polyglossia}
\setmainlanguage{arabic}
\setotherlanguage{english}
\newfontfamily\arabicfont[Script=Arabic]{Amiri}
\usepackage{geometry}
\geometry{margin=2.1cm}
\usepackage{hyperref}

\title{محاولة توليد Pair Kernel من آلية نقطية محلية متعددة المقاييس}
\author{مشروع تجربة البحث عن الأعداد الأولية}
\date{2026-09-01}

\begin{document}
\maketitle

\section{السؤال}

في التجربة السابقة استطعنا بناء نموذج نقطي يطابق التشتت ومنحنى
\[
\rho(k)
\]
عندما أضفنا حقلًا Gaussian ذا covariance مأخوذة مباشرة من نظرية الفترات القصيرة.

لكن ذلك كان bridge model، لا تفسيرًا نقطيًا مستقلًا.

السؤال الحالي:

\begin{quote}
هل يمكن أن تظهر covariance الزوجية المطلوبة من عمليات محلية على خط الأعداد نفسه،
بدل إدخال مصفوفة covariance جاهزة على مستوى الفترات؟
\end{quote}

لا نستخدم أي معلومات ثلاثية \(k=3\).

\section{الفكرة}

نستعمل random-residue sieve نفسه:
\[
q_0=7,\qquad Q=7919.
\]

ثم، بدل إضافة ضوضاء مستقلة، نسمح بعمليات إعادة توزيع محلية ذات مجموع صفري.
الأداة هي موجة Haar محلية.

على قطعة طولها \(L\)، تأخذ الموجة القيمة:
\[
+1
\]
على النصف الأول، و:
\[
-1
\]
على النصف الثاني، أو العكس عشوائيًا.

إذن كل حدث محلي يقول بصورة مبسطة:

\begin{quote}
زد الميل إلى اختيار نقاط في نصف محلي، وخفضه بالمقدار نفسه تقريبًا في النصف المجاور.
\end{quote}

هذه الآلية تولد تغايرًا سلبيًا بين المناطق المتجاورة من غير إدخال
covariance matrix للفترات.

\section{لماذا عدة مقاييس؟}

نظرية Leung تعطي شكلًا يمتد عبر عدة مسافات:
\[
\rho(k)\sim-\Delta(k).
\]

إذن آلية ذات طول واحد فقط لن تكون كافية غالبًا.
استخدمنا مقاييس ثنائية:
\[
L=2,4,8,\ldots,32768
\]
خلية، وحجم الخلية:
\[
1024
\]
عددًا صحيحًا.

أي أن المقاييس الفيزيائية تمتد من بضعة آلاف إلى عشرات الملايين.

\section{كيف اخترنا أوزان المقاييس؟}

هذه نقطة منهجية مهمة.

لم نستخدم بيانات الأوليات لاختيار أوزان المقاييس.
بدل ذلك حسبنا، لكل مقياس Haar، covariance التي يولدها هندسيًا بين فتراتنا.

ثم حللنا مسألة non-negative least squares:
\[
\min_{v_s\ge0}
\left\|
\sum_s v_s C_s-C_{\rm theory}
\right\|^2,
\]
حيث \(C_{\rm theory}\) هو kernel المستخرج من short-interval pair theory.

إذن:

\[
\boxed{
\text{شكل الآلية دُرّب على النظرية، لا على سلسلة الأوليات المرصودة.}
}
\]

هذا لا يجعل النموذج اشتقاقًا من المبادئ الأولى، لكنه يختلف جوهريًا عن إدخال
مصفوفة covariance على مستوى الفترات: هنا covariance تظهر من عمليات محلية
ذات دعم محدود على خط الأعداد.

\section{نتيجة ملاءمة kernel النظري}

كانت ملاءمة مكتبة Haar للنظرية قوية:

\[
\mathrm{RMSE}_{\rm kernel}\approx0.00139.
\]

وبالنسبة إلى الشكل المطبع مقابل:
\[
\frac{\Delta(k)}{\Delta(1)},
\]
كان:
\[
\mathrm{RMSE}_{\rm shape}\approx0.0145.
\]

الأوزان غير الصفرية تركزت عند المقاييس:

\begin{center}
\begin{tabular}{rr}
\toprule
المقياس بالأعداد الصحيحة & وزن التباين\\
\midrule
65536   & 3.416730\\
131072  & 0.154359\\
262144  & 0.213415\\
524288  & 0.087955\\
1048576 & 0.028193\\
2097152 & 0.013561\\
\bottomrule
\end{tabular}
\end{center}

إذن أفضل تمثيل محلي للنظرية في هذا النطاق ليس nearest-neighbor فقط،
بل عملية متعددة المقاييس، وأقوى وزن ظهر قرب:
\[
2^{16}=65536.
\]

يجب عدم تفسير هذا الرقم كقانون عالمي؛ قد يتحرك عندما نغير \(X\) أو أطوال الفترات.

\section{تقسيم التدريب والاختبار}

ولدنا \(40\) realizations مستقلة من random sieve.

استعملنا أول:
\[
12
\]
فقط لمعايرة معامل شدة واحد.

واستعملنا البقية:
\[
28
\]
للتقييم النهائي.

لكل realization في الاختبار ولدنا \(24\) حقول محلية مستقلة، أي:
\[
672
\]
محاكاة نهائية.

\section{المعامل الوحيد الذي أخذناه من البيانات}

بعد أن ثبتنا أوزان المقاييس من النظرية، بقي معامل شدة واحد:
\[
a.
\]

يعدل احتمال اختيار ناجٍ محليًا:
\[
p(x)
=
\theta_i+aF(x),
\]
حيث \(F(x)\) هو الحقل المحلي متعدد المقاييس.

حددنا \(a\) باستخدام بيانات التدريب فقط بحيث يطابق مقدار التشتت:
\[
R.
\]

حصلنا على:
\[
\boxed{
a\approx0.003833.
}
\]

في جميع محاكاة الاختبار بقي:
\[
0<p(x)<1,
\]
ولم نحتاج إلى clipping:
\[
\text{clip fraction}=0.
\]

\section{نتيجة الاختبار خارج العينة}

على \(672\) محاكاة اختبار:

\[
R_{\rm obs}=0.380332,
\]
و:
\[
R_{\rm model}=0.383104,
\]
مع:
\[
p\approx0.918.
\]

إذن معامل الشدة المعاير على مجموعة التدريب انتقل جيدًا إلى مجموعة الاختبار.

\section{الارتباط الأول}

البيانات:
\[
\rho_1^{\rm obs}=-0.123691.
\]

نظرية الفترات القصيرة:
\[
\rho_1^{\rm theory}=-0.108753.
\]

النموذج المحلي:
\[
\boxed{
\rho_1^{\rm local}=-0.116681.
}
\]

والقيمة المرصودة ليست شاذة تحت النموذج:
\[
p\approx0.875.
\]

هذه نتيجة أقوى من نموذج independent thinning القديم:
\[
\rho_1^{\rm independent}\approx-0.0808.
\]

\section{منحنى \(\rho(k)\)}

النموذج المحلي أعطى:
\[
\rho_2\approx-0.03375,
\]
بينما النظرية:
\[
-0.04438.
\]

وبالنسبة إلى ملاءمة:
\[
\rho(k)\approx-a_\Delta\Delta(k),
\]
وجدنا:
\[
a_\Delta^{\rm local}\approx0.15931,
\]
بينما النظرية لفتراتنا:
\[
a_\Delta^{\rm theory}\approx0.15871.
\]

أي أن مقدار الـamplitude الكلي ظهر تقريبًا صحيحًا من الآلية المحلية.

لكن shape RMSE بعد الانتقال من kernel Haar النظري إلى عملية الاختيار النقطية أصبح:
\[
\mathrm{RMSE}_{\rm shape}\approx0.0689.
\]

إذن الآلية لا تعيد الشكل النظري بنفس دقة bridge model:
\[
0.0370,
\]
لكنها تحقق جزءًا كبيرًا من الهدف من دون covariance جاهزة على مستوى الفترات.

\section{اختبار المنحنى الكامل}

اختبرنا:
\[
k=1,\ldots,20.
\]

إحصاء RMS المعياري للمنحنى أعطى:
\[
p_{\rm RMS}\approx0.358.
\]

واختبار أكبر انحراف مطلق:
\[
p_{\max}\approx0.186.
\]

إذن:
\[
\boxed{
\text{منحنى الارتباط المرصود ككل غير شاذ تحت الآلية المحلية.}
}
\]

lags 18 و20 منفردان يبدوان متطرفين، لكن الاختبار العالمي لا يرفض النموذج.

\section{اختبارات لم تدخل في المعايرة}

لم نستخدم runs أو triples لتحديد أوزان المقاييس أو \(a\).

مع ذلك:

\[
\mathrm{runs}_{\rm obs}=256,
\]
و:
\[
\mathrm{runs}_{\rm model}=242.93\pm10.82,
\]
مع:
\[
p\approx0.248.
\]

والثلاثيات الرتيبة:

\[
145
\]
مرصودة، مقابل:
\[
143.03
\]
في النموذج، مع:
\[
p\approx0.883.
\]

كما أن أنماط الإشارة الثمانية:
\[
000,\ldots,111
\]
لم تعط أي فشل إحصائي؛ جميع قيم \(p\) كانت أكبر من نحو:
\[
0.29.
\]

مرة أخرى لا يظهر حتى الآن سبب تجريبي لإضافة \(k=3\).

\section{المقارنة مع المراحل السابقة}

\begin{center}
\begin{tabular}{lrrrr}
\toprule
النموذج & \(R\) & \(\rho_1\) & \(\rho_2\) & runs\\
\midrule
البيانات & 0.3803 & -0.1237 & -0.0818 & 256.0\\
independent thinning & 0.4105 & -0.0808 & -0.0270 & 237.1\\
raw sieve/quota & 0.3085 & -0.1068 & -0.0360 & 240.9\\
pair-kernel bridge & 0.3810 & -0.1073 & -0.0375 & 241.0\\
local multiscale & 0.3831 & -0.1167 & -0.0338 & 242.9\\
\bottomrule
\end{tabular}
\end{center}

النموذج المحلي أفضل بوضوح من independent thinning في الحفاظ على
الارتباط، ويطابق التشتت في الوقت نفسه.

\section{ما الذي نجح؟}

نجحنا في نقل pair correction من:

\[
\text{Gaussian covariance على مستوى الفترات}
\]

إلى:

\[
\boxed{
\text{مجموع عمليات إعادة توزيع نقطية محلية متعددة المقاييس.}
}
\]

والآلية تستخدم فقط:

\begin{enumerate}
\item random sieve؛
\item موجات Haar محلية صفرية المجموع؛
\item أوزان مقاييس مأخوذة من kernel النظري؛
\item معامل شدة واحد مأخوذ من \(R\) على training set.
\end{enumerate}

\section{ما الذي لم ننجح فيه بعد؟}

الأوزان:
\[
v_s
\]
لم تُشتق بعد من قانون حسابي أو آلية أولية طبيعية.
لقد اخترناها بواسطة NNLS لكي تعيد \(\Delta(k)\).

لذلك لم نصل بعد إلى:

\begin{quote}
``هذه القاعدة المحلية البسيطة وحدها تولد تلقائيًا قانون Leung.''
\end{quote}

وصلنا إلى نتيجة أضعف لكنها مفيدة:

\begin{quote}
يمكن تمثيل kernel النظري بدقة عالية كمزيج من عمليات balancing محلية
محدودة الدعم، ويمكن تحويل هذا التمثيل إلى نموذج نقطي يعمل إحصائيًا.
\end{quote}

\section{السؤال التالي}

السؤال التالي أصبح:

\begin{quote}
هل للأوزان متعددة المقاييس قانون بسيط عندما نغير حجم \(X\) وطول الفترات؟
\end{quote}

أي بدل إعادة ملاءمة:
\[
v_s
\]
لكل تجربة، نقسم المجال إلى مستويات مختلفة من \(X\)، ونرى هل الأوزان تنهار
على منحنى scaling واحد بعد تطبيع المقياس بالنسبة إلى طول الفترة \(H\).

إذا ظهر قانون ثابت مثل:
\[
v(L/H)\sim f(L/H),
\]
فسيكون لدينا مرشح لآلية عامة، لا مجرد decomposition عددي لـ\(\Delta(k)\).

\section{ملفات التجربة}

\begin{itemize}
\item \texttt{local\_multiscale\_scale\_weights\_q7.csv}
\item \texttt{local\_multiscale\_theory\_kernel\_fit\_q7.csv}
\item \texttt{local\_multiscale\_point\_model\_rho\_q7.csv}
\item \texttt{local\_multiscale\_point\_model\_global\_rho\_test\_q7.csv}
\item \texttt{local\_multiscale\_point\_model\_sign\_triples\_q7.csv}
\item \texttt{local\_multiscale\_point\_model\_summary\_q7.csv}
\item \texttt{local\_multiscale\_model\_comparison\_q7.csv}
\item \texttt{analyze\_local\_multiscale\_pair\_mechanism.py}
\item \texttt{plots/local\_multiscale\_scale\_weights.svg}
\item \texttt{plots/local\_multiscale\_full\_rho\_comparison.svg}
\item \texttt{plots/local\_multiscale\_delta\_shape\_comparison.svg}
\end{itemize}

\end{document}
